\section{Proyecto eTwinning: «Sombras sin fronteras»}\label{sec:etwinning}

\subsection{Descripción}
eTwinning es una comunidad segura y gratuita en la que docentes y alumnado de centros europeos trabajan en proyectos colaborativos \parencite{esep}. En este proyecto, el IES Ejemplo y un centro de secundaria situado en torno a 50°~N miden en la misma semana, el mismo día si el cielo lo permite, la sombra de una varilla al mediodía solar y comparan cómo afecta la latitud a la altura del sol y a la separación de los paneles. Este planteamiento se adapta de la actividad escolar sobre Eratóstenes \parencite{eratostenes}.

\subsection{Objetivos}
\begin{itemize}
  \item Comprobar que el modelo matemático funciona en otro lugar y en otras condiciones (CE2, CE6).
  \item Comunicar datos, métodos y conclusiones a un público externo (CE8).
  \item Colaborar con otro centro de forma segura y respetuosa (CE9).
\end{itemize}

\subsection{Socio y puesta en marcha}
\begin{itemize}
  \item \textbf{Socio tipo:} grupo de edad similar en un centro de otro país eTwinning y otra latitud. La lengua de trabajo se acuerda entre ambos docentes.
  \item \textbf{Registro:} cada docente se registra con EU Login y solicita la validación de su organización nacional de apoyo \parencite{esep}.
  \item \textbf{Creación:} se necesitan dos socios fundadores; el proyecto debe ser aprobado por los Servicios Nacionales de Apoyo de ambos países \parencite{intef}. Se solicita al inicio del trimestre para que esté activo antes de S1.
\end{itemize}

\subsection{Calendario integrado en las sesiones}
\begin{tablaSA}{colspec={Q[c,wd=30pt]X[l]X[l]}, cells={font=\footnotesize}}
  Ses. & IES Ejemplo & Centro socio \\
  Antes & Registro, acuerdo de plan y calendario, apertura del TwinSpace & Igual \\
  1 & Presentación breve de cada equipo en el TwinSpace, con avatar & Presentación equivalente \\
  14 & Sombra de la varilla al mediodía solar; datos y fecha en una hoja compartida & Medición el mismo día, a su mediodía solar; con nubes, el primer día despejado de esa semana \\
  16 & Cálculo de $h$ de ambos centros; comparación de sombras y filas & Cálculo equivalente \\
  20 & Videomensaje de conclusiones y foro de preguntas & Respuesta y conclusiones conjuntas \\
\end{tablaSA}

\subsection{Resultado matemático esperado}
La altura del sol al mediodía del solsticio de invierno es $h=90^\circ-\varphi-23{,}45^\circ$. En Madrid ($40{,}4^\circ$) resulta $26{,}15^\circ$ y la sombra equivale a $2{,}04$ veces la altura del objeto; en $50^\circ$~N resulta $16{,}55^\circ$ y la sombra es $3{,}37$ veces la altura. El alumnado concluye que la misma instalación necesita más separación entre filas en el norte. El día de la medición, por ejemplo el 10 de marzo ($\delta\approx-4{,}9^\circ$), el sol está a unos $44{,}7^\circ$ en Madrid y a $35{,}1^\circ$ en $50^\circ$~N, y la varilla de 1~m da sombras de $1{,}01$~m y $1{,}42$~m. Si los centros miden en días distintos, cada altura se calcula con la declinación de su fecha. Como ampliación, con la diferencia de ángulos y la distancia entre centros sobre el mismo meridiano se estima el radio de la Tierra.

\subsection{Herramientas, seguridad y evaluación}
\begin{itemize}
  \item \textbf{Herramientas:} TwinSpace para foros y materiales, hoja de cálculo compartida, GeoGebra y videomensaje.
  \item \textbf{Seguridad y datos:} se usan avatares y nombres de pila o equipos, sin fotografías identificables del alumnado, y las autorizaciones familiares que exija el centro. El profesorado supervisa la comunicación.
  \item \textbf{Evaluación:} la comparación con el socio forma parte del libro (R1, criterios 6.1 y 2.2). El proyecto se valora con una encuesta conjunta a estudiantes y docentes sobre colaboración, aprendizaje y organización.
\end{itemize}
